\documentclass[12pt]{article}
\usepackage[utf8]{inputenc}
\usepackage[T1]{fontenc}
\usepackage{hyperref}
\title{Kiunum Language Resources for K-12 Curriculum Development}
\author{Compiled by Claude Sonnet 5 (claude-sonnet-5) for Andr\'as Kornai}
\date{}
\begin{document}
\maketitle
\section{Introduction}
Kiunum (ISO 639-3 wei) is also listed as Were or Weredai, and under the village name Lewada-Dewara. Search results place it in Dewara village, Gogodala Rural LLG, Western Province, Papua New Guinea. Sources give about 700 speakers, against the roughly 490 in the brief. Its classification is unsettled: Wikipedia's infobox says Trans-New Guinea, while the Anim-family literature appears to place a Lewada-Dewara language in the Tirio (Lower Fly) group. The search found very little about this language.
\section{Primary Educational Resources}
None found. No primers, textbooks or school programmes for Kiunum appeared in search results.
\section{Linguistic and Reference Materials}
\begin{itemize}
\item Wikipedia entry for the language (ISO 639:wei): \url{https://wikipedia.com/wiki/ISO_639:wei}
\item Numeral data, Max Planck database (traditional numerals 1--5 only): \url{https://lingweb.eva.mpg.de/channumerals/Kiunum.htm}
\item Usher and Suter, ``The Anim Languages of Southern New Guinea'': \url{https://muse.jhu.edu/article/586710/}
\item Anim languages overview: \url{https://en.wikipedia.org/wiki/Anim_languages}
\end{itemize}
\section{Religious and Community Resources}
None found.
\section{Digital and Audio Resources}
None found.
\section{Curriculum Development Context}
There is no evidence of a standard orthography, published teaching materials or school use. A curriculum would first need documentation work with the community, such as a word list, an orthography and recorded texts. The sources found are academic or encyclopedic only.
\end{document}